% Transcription — fonds Alexandre Grothendieck, shelfmark 12, batch 2 (pages 21-40).
% Established from the facsimile archives/batches/12/batch-02.pdf (a mirror of
% https://grothendieck.umontpellier.fr/12.pdf), per the skill
% transcribe-grothendieck. The batch file carries no cover sheet: PDF page N
% is page 20+N of the folder (the Montpellier footer prints « 21 » ... « 40 »),
% and the archivists' pencil numbers bottom-left agree leaf by leaf (21, 31,
% 32, 33, 34, 35, 36, 37, 38, 39, 40 are legible). He does not paginate this
% run himself.
%
% Pass: Opus 5.5 (claude-opus-5-5), 2026-09-23 — first pass, unchecked against the pages by a human.
%
% Nineteen of the twenty pages are transcribed. Page 21 is a leaf of an
% unrelated typescript (a carbon on generic flatness, Ass F, « prop. 3 a)
% (ii) », struck through by one long stroke, nothing in his hand) and is
% absent.
%
% What the batch is. One continuous run in his hand, headed « Motifs » on
% page 22, written out fairly and quickly, in numbered paragraphs 1) to 10)
% which follow, one for one, the first ten points of the table « Motifs »
% that closes batch 1 (page 20); batch 3 opens on 11). Paragraphs 9) and 10)
% carry underlined headings of his and become \section here, as batch 3
% does for 11) to 15); 1) to 8) have no headings.
%   22     1) M^+(X), effective motives over a noetherian prescheme: an
%          abelian Q-category with ⊗ (right exact), unit 1_X = Q_X(0);
%          D^b(M^+(X)) with ⊗^L. 2) f^* exact, Lf^*, compatible with ⊗;
%          Rf_* when f is of finite type and proper or Y excellent.
%   23     transitivity, projection formula; margin: the whole formalism
%          Rf_*, f^*, f^!, Rf_!. 3) X = lim X_i with affine transition maps:
%          M^+(X) ≃ lim M^+(X_i), whence reduction to finite type over Z.
%   24     the same for Rf_*. 4) T_ℓ : M^+(X) → M_ℓ(X), ℓ-adic realisation
%          into constructible Q_ℓ-sheaves: exact, faithful, not fully
%          faithful, compatible with ⊗, units, f^*, Rf_*.
%   25     extension to D^b; a cancelled block on T_∞ for Q-preschemes;
%          Künneth; 5) the canonical object Q_X(-1).
%   26     T_ℓ(Q(-1)) ≃ Q_ℓ(-1); Q(-1) = R^2 f_*(1) for P^1_X; R^{2d}f_*(1_X)
%          ≃ Z_S(-d) for f smooth projective with connected fibres, and
%          R^{2d}f_!(1_X) for f quasi-projective.
%   27     purity Ri^!(1_X) ≃ Q_Y(-d); 6) ⊗Q(-1) fully faithful, not an
%          equivalence; M(X) as the pseudo-inductive limit under ⊗Q(-1).
%   28     every object is M_n(n); T_ℓ extends. 7) internal Hom, RHom,
%          compatible with T_ℓ.
%   29     8) independence of ℓ: « constant tordu » for T_ℓ(M) and T_ℓ'(M)
%          together, same rank; traces and characteristic polynomials of an
%          endomorphism u are rational and independent of ℓ (boxed).
%   30     proof via the contraction M^∨ ⊗ M → 1_X and Hom(1_X, 1_X) = Q
%          for X connected (boxed). 9) filtrations: 1°) by weight,
%          M^{+0} ⊂ M^{+1} ⊂ …, checked fibre by fibre and at closed points.
%   31-33  local nature, Rf_! raises weight, generation over an algebraically
%          closed field by subquotients of R^i f_!(Z·1_X); multiplicativity,
%          Z(-1) of weight 2; extension to M(X) by twists, ⊗Q(-1) : M^i ≃
%          M^{i+2}; the inclusion Q(-j) ⊗ M^{+i} ⊂ M^{+i+2j} is strict;
%          thickness; the graded pieces G_i^+(X), G_i(X), all equivalent to
%          G_0(X) by twisting.
%   34-35  b) the filtration by « type dimensionnel » D_i: the smallest thick
%          filtration above the weight filtration and stable under ⊗Q(-1);
%          behaviour under f^*, R^j f_!, R^j f_* (the last queried);
%          extension to M(X), D_i ⊗ D_j ⊂ D_{i+j}, Ext^i(D_j, D_k) ⊂ D_{j+k}.
%   35-40  10) « Motifs constants tordus. Anneaux M^+(X) et M(X) »: the
%          subcategory M^+(X)_spéc, abelian, finite length, a « catégorie
%          tensorielle sur Q »; its ring M^+(X) = Z^(Σ^+(X)) on the classes
%          of simple objects, ξ = cl Q_X(-1), M(X) = M^+(X)[1/ξ]; gradings,
%          M_0(X), M^pair, M^impair; D_i(Σ^+(X)); primitive elements σ_i(X)
%          and the decomposition Σ_i = σ_i ∪ ξσ_{i-2} ∪ …; the quotients
%          D_i/D_{i-1}; the multiplication maps σ_i × σ_j → Z^(…); the
%          operations ∨ (x ↦ ξ^{-i}x expected) and λ^i (in terms of finite
%          groups such as the symmetric groups). Page 40 ends the paragraph;
%          the lower half of the leaf is blank.
%
% Notation fixed for the batch, agreeing with batches 1 and 3:
% \mathcal{M}^{+}(X), \mathcal{M}(X) for his script M (categories);
% \mathcal{M}^{+i}(X), later \mathcal{M}^{+}_{i}(X) — he switches from upper
% to lower indices on page 33, after writing « indices en bas » in the
% margins of pages 30 to 33; plain M^{+}(X), M(X) for the rings (he writes an
% upright M from page 35 on). \mathcal{G}_i for his curly G (graded pieces),
% D_i for the dimensional filtration, \mathfrak{S}_p for the looped S of
% page 40. \mathbf{1}_X is his 1_X; \mathbb{Q}(-1) and \mathbb{Z}(-1) are
% both on the leaves and are kept as written. His « tt », « isom », « t.f. »,
% « gén. », « ens. », « noeth. » are kept.
%
% Legibility. Formulas carry; the connecting prose of pages 31 to 40 is fast
% and often read only from the argument; where it would not resolve it is
% \ill{}. The slanted margins of pages 22, 23, 25, 27, 31 and 34 were turned
% upright before reading. The right margin of page 34 did not resolve.
% Nothing on these leaves is dated.
%
% The dating is the inventory's own for the folder, copied verbatim. Its
% group, [10-18] « Théorie des motifs », is dated [à partir de 1958-à partir
% de 1982].
\documentclass[11pt,a4paper]{article}
% Le préambule commun à toutes les transcriptions du fonds Grothendieck.
%
% Il tient en une idée : **l'incertitude se compose, elle ne se résout pas.**
% Une transcription qui lisse un mot douteux a détruit la seule chose pour
% laquelle on la faisait — un lecteur doit pouvoir savoir, à chaque endroit,
% ce qui était sur le feuillet et ce qui a été ajouté. Les six macros qui
% suivent sont donc l'appareil critique, et non de la décoration.
%
% Les mêmes commandes sont reconnues par `scripts/render.mjs`, qui produit la
% vue de lecture du site. Une macro ajoutée ici doit l'être aussi là-bas,
% faute de quoi le rendu échouera bruyamment — ce qui est voulu : un rendu
% silencieusement faux est pire qu'un rendu absent.

\ProvidesPackage{grothendieck}[2026/08/08 Transcriptions du fonds Grothendieck]

\usepackage{amsmath,amssymb,amsthm}
\usepackage{mathrsfs}
\usepackage{stmaryrd}
% Les diagrammes commutatifs sont partout dans ce fonds : c'est la langue
% même de la géométrie algébrique de Grothendieck, et une transcription qui
% les rendrait en prose ne transcrirait rien.
\usepackage{tikz-cd}
% Français par défaut — c'est la langue des feuillets — mais l'anglais est
% chargé aussi : la traduction et le résumé sont des documents anglais, et la
% typographie française (espaces avant les deux-points) y serait une faute.
% Ces documents basculent par \selectlanguage{english} en tête de corps.
\usepackage[english,french]{babel}
\usepackage{fontspec}
\usepackage{geometry}
\geometry{a4paper,margin=25mm,marginparwidth=28mm}
\usepackage{marginnote}
\usepackage{xcolor}
% ulem, not soul. `soul`'s \st and \ul are built on a character-by-character
% scan that XeLaTeX's font machinery breaks: the document fails with an
% undefined control sequence rather than degrading. ulem does the same two jobs
% and is Unicode-engine safe. `normalem` stops it hijacking \emph, which the
% transcriptions use for Grothendieck's own emphasis.
\usepackage[normalem]{ulem}
\usepackage{hyperref}
\hypersetup{colorlinks=true,linkcolor=black,urlcolor=black,pdfborder={0 0 0}}

% --- Métadonnées du lot -----------------------------------------------------
% Reprises telles quelles par le script de rendu, qui les affiche en tête de
% la vue de lecture. Elles fixent ce que le fichier prétend couvrir : sans
% elles, une transcription flotte sans point d'attache dans un fonds de seize
% mille feuillets.

\newcommand{\folder}[1]{\gdef\grothFolder{#1}}
\newcommand{\batch}[1]{\gdef\grothBatch{#1}}
\newcommand{\leaves}[2]{\gdef\grothFirst{#1}\gdef\grothLast{#2}}
\newcommand{\foldertitle}[1]{\gdef\grothFoldertitle{#1}}
\gdef\grothFolder{?}\gdef\grothBatch{?}\gdef\grothFirst{?}\gdef\grothLast{?}\gdef\grothFoldertitle{}

% --- Le résumé --------------------------------------------------------------
%
% Il ouvre la lecture modernisée, sous le titre « Résumé », et il est composé
% en petit. Ce n'est pas un ornement : c'est le seul passage du document qui
% ne soit pas des mathématiques, et un lecteur doit pouvoir le reconnaître
% avant de l'avoir lu — puis le sauter s'il connaît déjà le sujet, ou s'y
% arrêter s'il ne le connaît pas. Le filet qui le suit marque la reprise du
% corps.

\newenvironment{resume}
  {\small\setlength{\parskip}{0.45em}}
  {\par\medskip\hrule\medskip}

% --- Mots-clés ---------------------------------------------------------------
%
% En anglais, en clôture du résumé de la lecture modernisée : le vocabulaire
% moderne sous lequel chercher ce que les feuillets construisent. Le manifeste
% du site (`npm run manifest`) les extrait pour étiqueter la cote entière —
% cette ligne est la source unique des tags, il n'y a pas de fichier à côté.
\newcommand{\keywords}[1]{\par\smallskip\noindent
  {\footnotesize\textsc{Keywords} --- \textit{#1}}\par}

% --- L'appareil critique ----------------------------------------------------

% Le numéro de feuillet, en marge. C'est l'ancre qui permet de régler un
% désaccord sur une formule en regardant *un* feuillet plutôt qu'une chemise
% de six cents. Le site s'en sert aussi pour tourner les pages du fac-similé
% quand on fait défiler la transcription.
\newcommand{\leaf}[1]{%
  \marginnote{\footnotesize\color{gray!70}\textsf{#1}}%
  \label{leaf:#1}\ignorespaces}

% Illisible : on ne devine pas. Un mot inventé qui se lit comme les autres est
% le pire résultat possible.
\newcommand{\ill}{\textcolor{red!60!black}{[\,\ldots\,]}}

% Lecture douteuse : le mot est donné, et signalé comme incertain.
\newcommand{\uncertain}[1]{\uline{#1}}

% Ajout de l'éditeur — une lettre manquante, un mot sous-entendu.
\newcommand{\add}[1]{\textcolor{blue!50!black}{[#1]}}

% Biffé par Grothendieck lui-même. Ce qu'il a rayé fait partie du document :
% une rature dit souvent où la pensée a changé de direction.
\newcommand{\struck}[1]{\sout{#1}}

% Note du transcripteur, sur l'état matériel du feuillet.
\newcommand{\note}[1]{\par\smallskip{\small\color{gray!60!black}\textit{[#1]}}\par\smallskip}

% Note marginale de Grothendieck, distincte du corps du texte.
\newcommand{\marginal}[1]{\par\smallskip\noindent\hspace{1em}%
  {\small\color{gray!60!black}#1}\par\smallskip}

% --- Titre ------------------------------------------------------------------

\AtBeginDocument{%
  \begin{center}
    {\large\bfseries Fonds Alexandre Grothendieck --- cote n\textsuperscript{o}~\grothFolder}\\[2pt]
    {\itshape\grothFoldertitle}\\[4pt]
    {\small feuillets \grothFirst--\grothLast\ (lot \grothBatch)}\\[2pt]
    {\footnotesize\color{gray!60!black}Transcription établie d'après le fac-similé numérisé
     par l'Université de Montpellier.\\
     Les lectures incertaines sont soulignées, les illisibles notées [\,\ldots\,],
     les ajouts entre crochets.}
  \end{center}
  \vspace{1em}}

\endinput


\folder{12}
\batch{2}
\pages{21}{40}
\dating{1964-[vers 1971]}
\watermark{Édition de démonstration}
\foldertitle{Généralités (sous-groupes de Galois motiviques) (1964 ?) : notes manuscrites (s.d.), tapuscrit (s.d.)}

\begin{document}

\section{Motifs}
\note{titre de sa main, souligné, en tête du feuillet 22. La rédaction qui
suit n'est pas paginée par lui ; ses paragraphes numérotés 1) à 10)
reprennent un à un les dix premiers points de la table « Motifs » du feuillet
20 (lot 1)}

\page{22}
1) À tout préschéma noeth.\ (éventuellement de type fini sur un anneau
noethérien) $X$ est associée une \emph{catégorie abélienne}
$\mathcal{M}^{+}(X)$, dite catégorie des motifs effectifs sur $X$. C'est une
$\mathbb{Q}$-catégorie \struck{de $\mathbb{Q}$-modules} abélienne, i.e.\ pour
tt $M \in \operatorname{Ob} \mathcal{M}^{+}(X)$, et $n \in \mathbb{Z}$,
$n \neq 0$, $n\, \mathrm{id}_{M}$ est un isom de $M$. De plus,
$\mathcal{M}^{+}(X)$ est muni d'un produit tensoriel commutatif et unitaire
(exact à droite), l'unité est notée $\mathbf{1}_{X}$ ou $\mathbb{Q}_{X}(0)$.
\marginal{On peut \uncertain{en tirer} \ill{} les constructions des foncteurs
$\wedge^{i} M$, etc.…}

On considère aussi la catégorie dérivée bornée $\mathrm{D}^{b}(\mathcal{M}^{+}(X))$
de $\mathcal{M}^{+}(X)$. Le produit tensoriel est étendu en un bifoncteur
$M \overset{L}{\otimes} N$ pour $M, N \in \mathrm{D}^{b}(\mathcal{M}^{+}(X))$.

2) Soit $f : X \to Y$ un morphisme de préschémas \struck{\ill{}} noeth., il
lui est associé un \emph{foncteur exact}
\[
  f^{*} : \mathcal{M}^{+}(Y) \to \mathcal{M}^{+}(X),
\]
d'où
\[
  Lf^{*} : \mathrm{D}^{b}(\mathcal{M}^{+}(Y)) \to \mathrm{D}^{b}(\mathcal{M}^{+}(X)),
\]
compatible avec $\otimes$. \marginal{(\uncertain{Cond.\ de} transitivité.)}

Si $f$ est de type fini, et propre \emph{ou} $Y$ excellent, on a
\struck{défini} de même un foncteur
\[
  Rf_{*} : \mathrm{D}^{b}(\mathcal{M}^{+}(X)) \to \mathrm{D}^{b}(\mathcal{M}^{+}(Y))
\]

\page{23}
\marginal{Considérer aussi le \uncertain{formulaire} de dualité entre $Rf_{*}$,
$f^{*}$, et les foncteurs $f^{!}$, $Rf_{!}$ et leurs relations,
\struck{de dualité} et enfin les formules \uncertain{d'Alexander} reliant tous
ces foncteurs —}
satisfaisant une formule de transitivité, une formule de projection
\[
  Rf_{*}(M \otimes Lf^{*}(N)) \simeq Rf_{*}(M) \otimes N .
\]

3) Supposons $X = \varprojlim X_{i}$, système projectif filtrant à morph.\ affines.
Alors pour les foncteurs images inverses, on a
\[
  \mathcal{M}^{+}(X) \simeq \varinjlim \mathcal{M}^{+}(X_{i}) .
\]
En particulier, si $X$ est de type fini sur $S = \operatorname{Spec}(A)$,
alors $X$ est limite de préschémas $X_{i}$ de type fini sur $\mathbb{Z}$, et
la détermination de $\mathcal{M}^{+}(X)$, avec ses structures déjà envisagées
est \struck{dit} ramenée au cas des préschémas de type fini.

De même, si $(S_{i})$ est un système projectif filtrant à morph.\ affines,
$S = \varprojlim S_{i}$, et si $X$, $Y$ \ill{} de t.f.\ sur $S$, \ill{}
définis par $X_{i}$, $(Y_{i})$ de la façon habituelle, si

\page{24}
$M_{i_{0}} \in \operatorname{Ob} \mathrm{D}^{b}(\mathcal{M}^{+}(X_{i_{0}}))$,
d'où $M_{i}$, $M$, on a pour $f_{i} : X_{i} \to Y_{i}$ la relation
\[
  Rf_{*}(M) = \varinjlim v_{i}^{*}(Rf_{i*}(M_{i}))
\]
où $v_{i} : Y \to Y_{i}$ est le morphisme canonique.

4) Soit $\ell \in \mathbb{P}$ tel que $\ell \mathbf{1}_{X} \in \Gamma(X,
\mathcal{O}_{X})$ soit inversible. Alors on a un foncteur
\[
  T_{\ell} = T_{\ell}^{(X)} : \mathcal{M}^{+}(X) \to \mathcal{M}_{\ell}(X)
\]
où $\mathcal{M}_{\ell}(X)$ est la catégorie des « $\mathbb{Q}_{\ell}$-modules
constructibles sur $X$ », i.e.\ la catégorie déduite de la catégorie des
« systèmes projectifs $\ell$-adiques de faisceaux de $\ell$-torsion
constructibles » en négligeant \ill{} les faisceaux de torsion. Le foncteur
$T_{\ell}$ est compatible avec $\otimes$ et \emph{unités}, \emph{exact} et
\emph{fidèle} (mais non pleinement fidèle), \emph{compatible} avec \emph{les}

\page{25}
\marginal{compatibilité avec $f^{!}$, $Rf_{!}$, avec \ill{} etc.…}
\emph{changement de base} $f^{*}$, et \emph{compatible} également \emph{avec}
$Rf_{*}$. [N.B. $T_{\ell}$ s'étend évidemment en un foncteur
$\mathrm{D}^{b}(\mathcal{M}(X)) \to \mathrm{D}^{b}(\mathcal{M}_{\ell}(X))$.]
$\{$La détermination des $T_{\ell}$ est encore ramenée au cas où $X$ est de
t.f.\ sur $\mathbb{Z}$.

Si $X$ est un $\mathbb{Q}$-préschéma de type fini, on a aussi un foncteur
\[
  T_{\infty} = T_{\infty}^{(X)} : \mathrm{D}^{b}(\mathcal{M}(X)) \to \mathcal{M}_{\infty}(X)
\]
où $\mathcal{M}_{\infty}(X)$ est la catégorie des Modules \ill{}
\uncertain{cohérents}. Ce foncteur $T_{\infty}$ est encore exact et
\emph{fidèle}, compatible avec $f^{*}$ et $Rf_{*}$.
\note{tout ce paragraphe sur $T_{\infty}$ est encadré puis biffé de deux
diagonales ; « de type fini » est ajouté au-dessus d'un mot biffé}
\marginal{NB. Ceci exclut déjà la \uncertain{chose} triviale
$\mathcal{M}^{+}(X) = 0$ pour tt $X$, car il faudrait que
$R^{*}f_{*}(\mathbb{Z}_{\ell}) = 0$ pour $f$, $\ell$, \ill{} ce qui n'est pas
vrai en général…}

Signalons aussi la compatibilité de $T_{\ell}$ avec l'isom.\ de
\emph{Künneth}.

5) Pour tout $X$, on a un \struck{\ill{}} élément canonique
\[
  \mathbb{Q}_{X}(-1) \quad\text{ou}\quad \mathbf{1}_{X}(-1) \in \operatorname{Ob} \mathcal{M}^{+}(X)
\]

\page{26}
dont la formation est compatible avec changement de base (il suffit donc de
le considérer sur $\operatorname{Spec} \mathbb{Z}$), avec des isom
\[
  T_{\ell}(\mathbb{Q}(-1)) \simeq \mathbb{Q}_{\ell}(-1)
  = \underbrace{T_{\ell}(\mathbb{G}_{m})^{-1}}_{\text{Module de Tate inverse de } \mathbb{G}_{m}}
\]
et le cas échéant \ill{} ($X$ sur $\mathbb{Q}$). On peut définir
$\mathbb{Q}(-1)$ comme $R^{2} f_{*}(\struck{\mathbb{Z}}\, \mathbf{1}_{\mathbb{P}^{1}_{X}})$,
où $f : \mathbb{P}^{1}_{X} \to X$ est la projection canonique. Posant
$\mathbb{Q}(-n) = \mathbb{Q}(-1)^{\otimes n}$ pour $n \geqslant 0$ on peut
prouver, à l'aide des axiomes déjà posés, que si $f : X \to S$ est lisse
projectif à fibres géom.\ connexes et \uncertain{non vides}, partout de dim
relative $d$, alors
\[
  R^{2d} f_{*}(\mathbf{1}_{X}) \simeq \mathbb{Z}_{S}(-d)
\]
et si on oublie l'hyp « projectif » mais seulement $f$ quasi-projectif, on
trouve encore
\[
  R^{2d} f_{!}(\mathbf{1}_{X}) \simeq \mathbb{Z}_{S}(-d)
\]

\page{27}
On veut \struck{exige} de plus \struck{compatibilité}, que si $X/S$ est lisse
et $Y \overset{i}{\subset} X$ est lisse sur $S$, de codim $d$ dans $X$,
l'isom
\[
  Ri^{!}(\mathbf{1}_{X}) \simeq \mathbb{Q}_{Y}(-d)
\]
compatible avec les isom déjà connus du point de vue $\ell$-adique —.

6) \struck{Chaque $\mathcal{M}(X)$ est plongé dans} Le foncteur
$M \mapsto M(-1) = M \otimes \mathbb{Q}(-1)$ de $\mathcal{M}^{+}(X)$ dans
lui-même est \emph{pleinement} fidèle mais pas une équivalence en général. Il
y a donc une façon canonique \uncertain{d'élargir} $\mathcal{M}^{+}(X)$ en
$\mathcal{M}(X)$ de telle façon que $\otimes \mathbb{Q}(-1)$ devienne une
équivalence : en prenant la pseudo-limite inductive de \ill{} des
\[
  \mathcal{M}^{+} \xrightarrow{\otimes \mathbb{Q}(-1)} \mathcal{M}^{+}
  \xrightarrow{\otimes \mathbb{Q}(-1)} \mathcal{M}^{+} \to \cdots
\]
On prolonge à $\mathcal{M}(X)$ la structure $\otimes$, dont $\mathbb{Q}(-1)$
devient inversible, d'inverse $\mathbb{Q}(1)$ \ill{}, et tt élément de
$\mathcal{M}(X)$ peut s'écrire,
\marginal{$\mathcal{M}^{+}(X)$ est sous-catégorie pleine \uncertain{épaisse}
abélienne de $\mathcal{M}(X)$ ; \ill{}}

\page{28}
$n$ assez grand, sous la forme $M_{n}(n)$, avec
$M_{n} \in \operatorname{Ob} \mathcal{M}^{+}(X)$ [pour $n$ fixé, $M_{n}$ est
bien déterminé par $M$ à isom unique près, et
$(M_{n}(n) \simeq M_{n+1}(n+1)) \Leftrightarrow (M_{n+1} \simeq M_{n}(-1))$ ;
on retrouve la description de $\mathcal{M}(X)$ en termes de \ill{}
$\varinjlim$.]. $T_{\ell}$ s'étend à $\mathcal{M}(X)$, de façon unique à
isom unique près, compatible avec $\otimes$.

Les foncteurs \ill{}

7) Dans $\mathcal{M}(X)$, on a aussi des foncteurs
$\underline{\mathrm{Hom}}$, liés à $\otimes$ par la formule habituelle
\[
  \mathrm{Hom}(P \otimes Q, R) \simeq \mathrm{Hom}(P, \underline{\mathrm{Hom}}(Q, R))
  \simeq \mathrm{Hom}(Q, \underline{\mathrm{Hom}}(P, R))
\]
et \struck{satisfaisant} qui s'étendent en
\[
  R\underline{\mathrm{Hom}}(P, Q) \qquad P, Q \in \operatorname{Ob} \mathrm{D}^{b}\mathcal{M}(X)
\]
satisfaisant à la relation analogue relativement à $\overset{L}{\otimes}$.
La formation des $\underline{\mathrm{Hom}}$ et $R\underline{\mathrm{Hom}}$ est
\emph{compatible} avec les $T_{\ell}$.
\note{les soulignements des Hom de la formule ne se distinguent pas tous ;
on a souligné les Hom internes}

\struck{N.B.} [N.B. On retrouve la formation des $f^{!}$ —]

\page{29}
8) Soient $\ell, \ell' \in \mathbb{P}$ premiers aux car.\ rés.\ de $X$. Soit
$M \in \operatorname{Ob} \mathcal{M}^{+}(X)$. On veut que
\[
  \begin{array}{c}
    T_{\ell}(M) \text{ faisceau constant tordu} \\
    \Updownarrow \\
    T_{\ell'}(M) \text{ faisceau constant tordu}
  \end{array}
\]
et que sous ces conditions, $T_{\ell}(M)$ et $T_{\ell'}(M)$ doivent avoir
même rang en chaque pt.

Pour vérifier l'égalité des rangs, on est ramené au cas où $X$ est le spectre
d'un corps (qu'on peut supposer si on veut, alg.\ clos). Plus généralement,
si $u$ est un endomorphisme de $M$, et ($X = \operatorname{Spec} k$, $k$ un
corps), on en déduit
\[
  T_{\ell}(u) \in \operatorname{End} T_{\ell}(M), \qquad
  T_{\ell'}(u) \in \operatorname{End}(T_{\ell'}(M))
\]
et on désire que
\[
  \boxed{\operatorname{Tr} T_{\ell}(u) = \operatorname{Tr} T_{\ell'}(u) \in \mathbb{Q}}
\]
d'où, remplaçant $u$ par $\wedge^{n} u$, la formule
\[
  \boxed{P(T_{\ell}(u), t) = P(T_{\ell'}(u), t) \in \mathbb{Q}[t]}
\]

\page{30}
[où il s'agit des polynômes caractéristiques]. Pour ceci, notons que
\[
  u \in \operatorname{Hom}(\mathbf{1}_{X}, \underline{\mathrm{Hom}}(M, M))
  = \operatorname{Hom}(\mathbf{1}_{X}, M^{\vee} \otimes M)
\]
et on a un morphisme \emph{contraction}
\[
  M^{\vee} \otimes M \to \mathbf{1}_{X}
\]
\marginal{à inclure dans le formalisme tensoriel}
d'où un $c(u) \in \operatorname{Hom}(\mathbf{1}_{X}, \mathbf{1}_{X})$ et
\struck{\ill{}}
\[
  \operatorname{Tr} T_{\ell}(u) = T_{\ell}(c(u)) \in \mathbb{Q}_{\ell} .
\]
et il suffit de savoir :
\[
  \boxed{X \text{ connexe} \Longrightarrow \operatorname{Hom}(\mathbf{1}_{X}, \mathbf{1}_{X}) = \mathbb{Q}\, \mathrm{id}_{\mathbf{1}_{X}}}
\]

\section{9) Filtration de $\mathcal{M}^{+}(X)$ et $\mathcal{M}(X)$}
\note{titre de sa main, souligné, au milieu du feuillet 30}

\subsection{1°) Filtration par le poids}

\marginal{indices en bas}
\struck{\ill{}}
\[
  \mathcal{M}^{+0}(X) \subset \mathcal{M}^{+1}(X) \subset \cdots \subset \mathcal{M}^{+i}(X) \subset \cdots
\]
filtration \emph{exhaustive} de $\mathcal{M}^{+}(X)$. [L'appartenance à
$\mathcal{M}^{+i}(X)$ se vérifie fibre \uncertain{sur $\mathbb{Z}$} par fibre,
\struck{\ill{}} (géométriques si on veut) [et dans le cas où $X$ de t.f.] il
suffit de vérifier en ses pts fermés. (Pour \ill{}, il y a un critère par
Frobenius, cf.

\page{31}
\marginal{\uncertain{Généraliser} compatibilité avec $\varprojlim$ des
préschémas…}
plus bas).] \struck{$\mathcal{M}^{i}(X)$ est « engendré » par les $R^{i}f_{*}$}
\uncertain{loc.\ pour} \ill{} L'appartenance à $\mathcal{M}^{i}(X)$ \ill{} de
\struck{$M|x$ implique l'appartenance} \ill{} compatible \ill{} $f^{*}$, et
\struck{l'appart.} $M|x \in \operatorname{Ob} \mathcal{M}^{+i}(k(x))
\Longrightarrow \exists\, U \ni x$ voisinage de $\bar{x}$, tel que
$M|U \in \operatorname{Ob} \mathcal{M}^{+i}(U)$. Soit $f : X \to Y$, alors
\[
  R^{j}f_{!} : \mathcal{M}^{+i}(X) \to \mathcal{M}^{+i+j}(Y) ;
\]
\note{l'exposant de $R$ se lit mal ; le $j$ est lu d'après le but}
de plus, $\mathcal{M}^{+i}(k)$ ($k$ un corps alg.\ clos) est « engendré » par
les \uncertain{sous-quotients} de
\[
  R^{i}f_{!}(\mathbb{Z} \cdot \mathbf{1}_{X})
\]
pour $X$ si on veut projectif lisse de dim $\leqslant j$ sur $k$
\marginal{indices en bas}
[on sait \uncertain{que tout} $M \in \operatorname{Ob} \mathcal{M}^{i}(k)$ a une
filtration dont les facteurs sont isomorphes à des tels sous-quotients].
\note{« $\mathcal{M}^{+i}(k)$ » est écrit par-dessus un $X$ biffé}

\struck{$\mathcal{M}^{i}(X) \otimes \mathbb{Z}(-1)$}
\[
  \left\{
  \begin{array}{l}
    \mathcal{M}^{+i}(X) \otimes \mathcal{M}^{+k}(X) \subset \mathcal{M}^{+i+k}(X) \\
    \mathbb{Z}(-1) \in \operatorname{Ob} \mathcal{M}^{+2}(X)
  \end{array}
  \right.
\]
d'où
\[
  \mathcal{M}^{+i}(X) \otimes \mathbb{Z}(-j) \subset \mathcal{M}^{+i+2j}(X)
  \qquad \text{pour } j \geqslant 0
\]

\page{32}
on a mieux :
\[
  M \in \operatorname{Ob} \mathcal{M}^{+i}(X) \Longleftrightarrow
  M(-j) \in \operatorname{Ob} \mathcal{M}^{+i+2j}(X)
\]
De cette façon, la filtration de $\mathcal{M}^{+}(X)$ par les
$\mathcal{M}^{+i}(X)$ ($i \geqslant 0$) permet de prolonger en une filtration
de $\mathcal{M}(X)$ par des $\mathcal{M}^{i}(X)$ ($i \in \mathbb{Z}$) de telle
façon que pour $i \geqslant 0$, on ait la filtration déjà
\struck{$\mathcal{M}(X)$} envisagée, et pour $i$ quelconque, on définit
\[
  M \in \operatorname{Ob} \mathcal{M}^{i}(X) \Longleftrightarrow
  M(-j) \in \operatorname{Ob} \mathcal{M}^{i+2j}(X)
\]
si on prend $j$ assez grand tel que
$M(-j) \in \operatorname{Ob} \mathcal{M}^{+}(X)$. On notera que
\[
  \mathcal{M}^{+i}(X) = \mathcal{M}^{i}(X) \cap \mathcal{M}^{+}(X)
\]
si $i \geqslant 0$, et pour tt $i$ si on définit
\[
  \mathcal{M}^{+i}(X) = \{0\} \struck{\mathcal{M}^{+}(X)} \quad \text{si } i < 0
\]
\struck{\ill{} $\mathcal{M}^{+}(X)$ \ill{} dans \ill{} par $\mathcal{M}^{+i}(X)$
\uncertain{consiste} \ill{}}
\marginal{indices en bas}

On aura
\[
  \begin{array}{rcl}
    \otimes \mathbb{Q}(-1) & : & \mathcal{M}^{i}(X) \overset{\text{équiv.}}{\simeq} \mathcal{M}^{i+2}(X) \\
    \otimes \mathbb{Q}(j) & : & \mathcal{M}^{i}(X) \overset{\text{équiv.}}{\simeq} \mathcal{M}^{i-2j}(X)
  \end{array}
\]
\note{les deux $\mathbb{Q}$ sont écrits par-dessus une autre lettre, biffée ;
l'exposant de la seconde ligne se lit $i-2j$ d'après l'argument}

\page{33}
mais \ill{} par l'inclusion
\[
  \begin{array}{l}
    \mathbb{Q}(-1) \otimes \mathcal{M}^{+i}(X) \hookrightarrow \mathcal{M}^{+(i+2)}(X) \qquad \text{plus gén.} \\
    \mathbb{Q}(-j) \otimes \mathcal{M}^{+i}(X) \hookrightarrow \mathcal{M}^{+(i+2j)}(X)
  \end{array}
\]
est \emph{stricte} en général, i.e.\ n'est pas une équivalence de catégories.

\marginal{indices en bas}
Notons que $\mathcal{M}^{+i}(X)$ est \emph{épaisse} dans $\mathcal{M}^{+j}(X)$
\struck{[elle est définie dans la catégorie]}, et de même
\[
  \mathcal{M}^{i}(X) \text{ épaisse dans } \mathcal{M}^{j}(X) \qquad (i \leqslant j).
\]
Les quotients
$\mathcal{G}^{+}_{i}(X) = \mathrm{Gr}^{F}_{i}(\mathcal{M}^{+}(X)) \simeq
\mathcal{M}^{+}_{i}(X)/\mathcal{M}^{+}_{i-1}(X)$ et
$\mathcal{G}_{i}(X) = \mathrm{Gr}^{F}_{i}(\mathcal{M}(X)) =
\mathcal{M}_{i}(X)/\mathcal{M}_{i-1}(X)$ sont fort intéressants.

Notons que les $\mathcal{G}_{i}(X)$ sont \ill{}, \uncertain{équiv.\ entre
eux} \struck{isomorphes} par twisting $\mathbb{Z}(-j)$ :
\[
  \mathcal{G}_{i}(X) \simeq \mathcal{G}_{i+2j}(X) ,
\]
on \ill{} \uncertain{équiv.\ can.} \struck{\ill{}} à $\mathcal{G}_{0}(X)$.
D'ailleurs
\[
  \mathcal{G}^{+}_{i}(X) \subsetneq \mathcal{G}_{i}(X) \simeq \mathcal{G}_{0}(X)
\]
\note{un tilde est écrit sous l'inclusion stricte : il faut lire
« équivalente à une sous-catégorie », comme le dit l'ajout qui suit}
équiv.\ à une sous-catégorie pleine et \emph{épaisse} de
$\mathcal{G}_{i}(X) \simeq \mathcal{G}_{0}(X)$.

\page{34}
\subsection{b) Filtration par « le \struck{dimension} type dimensionnel »}

On pose pour $i \geqslant 0$
\[
  D_{i}(\mathcal{M}^{+}(X)) = \struck{\mathcal{M}^{i}} \text{ sous-catégorie pleine}
\]
de $\mathcal{M}^{+}(X)$ formée des objets qui se \uncertain{loc.\ étale} [(dans
\ill{}, \ill{} fibre par fibre)] \ill{} objets de la forme
\[
  \struck{\ill{}} \; M(-j), \quad \text{avec } M \in \operatorname{Ob} \mathcal{M}^{+i}(X), \; j \in \mathbb{Z},\ j \geqslant 0
\]
i.e.\ on prend (\ill{}) la filtration minimale qui majore celle par le poids,
\struck{\ill{}} qui soit stable par $\otimes \mathbb{Q}(-1)$, et
\emph{épaisse}.
\marginal{\ill{}}
\note{une note de la marge droite, écrite de bas en haut et en partie biffée,
ne se lit pas}

Les $D_{i}(\mathcal{M}^{+}(X))$ sont stables par l'image inverse, et par
l'image directe on a
\[
  \begin{array}{l}
    R^{j}f_{!}(D_{i}(\mathcal{M}^{+}(X))) \subset D_{i+j}(\mathcal{M}^{+}(Y)) \\
    R^{j}f_{*}(D_{i}(\mathcal{M}^{+}(X))) \subset D_{i+j}(\mathcal{M}^{+}(Y)) \quad \text{?}
  \end{array}
\]
\note{le point d'interrogation de la seconde ligne est cerclé dans la marge}
\marginal{\ill{} formellement les filtrations $D_{i}(\mathcal{M})$ et
$\mathcal{M}_{i}$ !}
\note{cette note de la marge gauche est reliée par une flèche à la ligne de
$R^{j}f_{*}$}
[\uncertain{On a seulement}
$R^{j}f_{*}(\mathcal{M}^{+(i)}(X)) \subset \mathcal{M}^{+(i+2j)}(Y)$ !?]

On \struck{\ill{}} a :
\[
  M \in \operatorname{Ob} D_{i}(\mathcal{M}^{+}(X)) \Longrightarrow
  M(-j) \in \operatorname{Ob} D_{i}(\mathcal{M}^{+}(X))
\]
\marginal{pas $D_{i+j}$ !!}
\note{le $D_i$ du second membre est écrit par-dessus un autre indice et
cerclé, la flèche menant à la note « pas $D_{i+j}$ !! »}
\marginal{\uncertain{tout se passe} \ill{} en termes de la dim.\ de $X$ !}
et la filtration de $\mathcal{M}^{+}(X)$ par les $D_{i}(\mathcal{M}^{+}(X))$
s'étend en une filtration de $\mathcal{M}(X)$ par des $D_{i}(\mathcal{M}(X))$,

\page{35}
\uncertain{en} utilisant la filtration donnée, \struck{\ill{}} i.e.\ on a
\[
  M \in \operatorname{Ob} D_{i}(\mathcal{M}(X)) \Longleftrightarrow
  M(-j) \in \operatorname{Ob} D_{i}(\mathcal{M}(X))
\]
pour $j$ \uncertain{qui} \ill{}.

On a
\struck{$D_{i}(\mathcal{M}^{+}(X))/D_{i-1}(\mathcal{M}^{+}(X)) \xrightarrow{\ \approx\ } D_{i}(\mathcal{M}(X))/D_{i-1}(\mathcal{M}(X))$}
\[
  D_{i}(\mathcal{M}(X)) \otimes D_{j}(\mathcal{M}(X)) \subset D_{i+j}(\mathcal{M}(X))
\]
(mais $\mathbb{Z}(-j) \in \operatorname{Ob} D_{0}(\mathcal{M}(X))$ pour tt $j$.)
et de $=$ en mettant des $+$.

Enfin, on a
\[
  \underline{\underline{\mathrm{Ext}}}^{i}(D_{j}(\mathcal{M}(X)), D_{k}(\mathcal{M}(X))) \subset D_{j+k}(\mathcal{M}(X))
\]
(par une formule aussi simple en termes des \uncertain{$\mathcal{M}_{\alpha}(X)$} !).

\section{10) Motifs constants tordus. Anneaux $M^{+}(X)$ et $M(X)$}
\note{titre de sa main, souligné en deux morceaux, au bas du feuillet 35}

Soit $\mathcal{M}^{+}(X)_{\mathrm{spec.}}$ la sous-catégorie de
$\mathcal{M}^{+}(X)$ formée des motifs effectifs \emph{constants tordus}.
[\uncertain{Par} condition locale, qui se vérifie sur $T_{\ell}(M)$,
cf.\ No 8], et de même on définit $\mathcal{M}(X)_{\mathrm{spec.}}$
\struck{N.A. par \ill{} $M$ et}. \uncertain{Ainsi}
\[
  M \text{ est tordu} \Longleftrightarrow M(j) \text{ est tordu.}
\]
$\mathcal{M}^{+}(X)_{\mathrm{spec.}}$ et $\mathcal{M}(X)_{\mathrm{spec.}}$
sont des sous-\emph{catégories abéliennes} \emph{stables par}
\uncertain{extensions}, \emph{dont tt objet est de}

\page{36}
\emph{longueur finie}. Elle est stable par $\otimes$, \ill{} cette catégorie
les $\mathrm{Ext}^{i}$ ($i > 0$) \uncertain{sont nuls}, et la catégorie est
aussi stable par $\underline{\mathrm{Hom}}$ (cas de
$\mathcal{M}(X)_{\mathrm{spec.}}$). C'est une « catégorie tensorielle sur
$\mathbb{Q}$ ».

Le \uncertain{$\lambda$-anneau} défini par cette catégorie est noté $M^{+}(X)$
resp.\ $M(X)$.

Notons que $\xi = \mathrm{cl}\, \struck{\mathbb{Z}}\mathbb{Q}_{X}(-1)$ est un
élément de $M^{+}(X)$, et on a
\[
  M(X) = M^{+}(X)_{\xi} = M^{+}(X)\bigl[\tfrac{1}{\xi}\bigr] .
\]
En fait, $\xi$ est non diviseur de $0$ : $M^{+}(X)$ est un $\mathbb{Z}$-module
libre admettant comme base les classes d'\struck{\ill{}} objets
\emph{simples} de $\mathcal{M}^{+}(X)$, et $x \mapsto \xi \cdot x$ dans
$M^{+}(X)$ est défini par une application \emph{injective} \ill{} des objets
simples de $\mathcal{M}^{+}(X)$. [N.B. $M$ simple $\Longleftrightarrow$ $M$
simple dans $\mathcal{M}(X)$ $\Longleftrightarrow$ $M(-1)$ simple, comme il
résulte \ill{} que $\mathcal{M}^{+}(X)$ est épaisse dans $\mathcal{M}(X)$ et
$\otimes \mathbb{Z}(-1)$ une autoéquivalence]

\page{37}
Si $\Sigma^{+}(X) =$ ensemble des classes d'objets simples de
$\mathcal{M}^{+}(X)$,
\[
  \Sigma(X) = \varinjlim_{\text{par } e \mapsto \xi e} \Sigma^{+}(X) ,
\]
on a \struck{$M(K)$} un diagramme \ill{} :

\begin{tikzcd}
  M^{+}(X) \arrow[d, hook] & = & \mathbb{Z}^{(\Sigma^{+}(X))} \arrow[d, hook] \\
  M(X) & = & \mathbb{Z}^{(\Sigma(X))}
\end{tikzcd}

Alors $\mathcal{M}_{i}(X)$ se définit en termes \ill{}
\[
  \Sigma^{+}_{(i)}(X) \quad \text{et} \quad \Sigma_{i}(X),
\]
d'une partie \ill{}, d'où
\[
  M^{+}_{i}(X) = \mathbb{Z}^{\Sigma^{+}_{(i)}(X)}
\]
ce qui donne une filtration croissante sur l'anneau $M^{+}(X)$, et une
filtration analogue sur $M(X)$.

En fait, posant
\[
  \Sigma^{+}_{i}(X) = \Sigma^{+}_{(i)}(X) - \Sigma^{+}_{(i-1)}(X)
\]
on aura \ill{}
\[
  \Sigma^{+}_{i}(X) \cdot \Sigma^{+}_{j}(X) \subset \mathbb{Z}^{\Sigma^{+}_{i+j}(X)}
\]
donc $M^{+}(X)$ devient un \emph{anneau}

\page{38}
\emph{gradué} [i.e.\ à degrés $\geqslant 0$], et $\xi$ est de degré 1. De cette
façon, $M(X)$ devient gradué [i.e.\ à degrés de tout signe, $\xi$ étant un
élément inversible de degré 1]. Notons que
\[
  M_{0}(X) \simeq \mathbb{Z}^{\Sigma_{0}(X)}
\]
avec
\[
  \Sigma_{0}(X) = \varinjlim_{\text{par } x \mapsto \xi x} \Sigma^{+}_{i}(X)
\]
et
\[
  \begin{array}{l}
    M^{\text{pair}}(X) = M_{0}(X)[\xi] \qquad (\xi \text{ de degré } 1) \\
    M^{\text{impair}}(X) = \underline{M}_{1}(X) \otimes_{M_{0}(X)} M_{0}(X)[\xi]
  \end{array}
\]
\note{la seconde ligne est surchargée de petits ajouts au-dessus du produit
tensoriel, qui ne se lisent pas ; le $M_1$ est souligné}
\uncertain{donc} la structure d'anneau gradué de $M(X)$ est déterminée en
termes de la structure d'anneau de $M_{0}(X)$, \struck{et} la structure de
Module $M_{1}(X)$ sur $M_{0}(X)$, et l'application de mult.
\[
  M_{1}(X) \otimes_{M_{0}(X)} M_{1}(X) \to M_{0}(X)
\]
définie par \ill{}.

D'autre part, les objets de $D_{i}(\mathcal{M}^{+}(X))$ sont
\struck{\ill{}} \uncertain{ceux} dont la classe \ill{} dans \struck{$M^{+}(X)$
sont} $M^{+}(X)$ sont dans
\[
  \mathbb{Z}^{D_{i}(\Sigma^{+}(X))}
\]
où
\[
  \begin{aligned}
    D_{i}(\Sigma^{+}(X)) &= \coprod_{j \geqslant 0} \xi^{j}\, \Sigma^{+}_{(i)}(X) \\
    &= \coprod_{\substack{j \geqslant 0 \\ 0 \leqslant k \leqslant i}} \xi^{j}\, \Sigma^{+}_{k}(X)
  \end{aligned}
\]
\note{un premier symbole après le signe $=$ est biffé ; à droite des deux
lignes, derrière une accolade, un ajout se lit à peu près
« $\xi^{j}\Sigma^{+}_{[j]}(X)$ », peut-être biffé}

\page{39}
Posons
\[
  \struck{\Sigma}\; \sigma_{i}(X) = \Sigma^{+}_{i}(X) - \bigcup_{\alpha \geqslant 0} \xi^{\alpha}\, \Sigma^{+}_{i-2\alpha}(X)
\]
\note{sic pour $\alpha \geqslant 0$ : pour $\alpha = 0$ la réunion contiendrait
$\Sigma^{+}_{i}(X)$ tout entier ; un indice biffé est écrit sous la réunion}
(donc on a une partition (au sens large))
\[
  \Sigma_{i}(X) = \sigma_{i}(X) \cup \xi\, \sigma_{i-2}(X) \cup \cdots \cup
  \left\{
  \begin{array}{l}
    \xi^{i/2}\, \sigma_{0}(X) \\
    \xi^{\frac{i-1}{2}}\, \sigma_{1}(X)
  \end{array}
  \right.
\]
\marginal{suivant la parité de $i$}
\marginal{\emph{éléments primitifs} de poids $i$}
\note{la note cerclée « éléments primitifs de poids $i$ » est reliée par un
trait à $\sigma_i(X)$ ; la note « suivant la parité de $i$ », à droite de
l'accolade, porte des traits qui la biffent peut-être}
et on a donc
\[
  D_{i}(\Sigma^{+}(X)) \cap \Sigma^{+}_{j}(X) =
  \begin{cases}
    \Sigma^{+}_{j}(X) & \text{si } j \leqslant i \\
    \xi^{j-i}\, \sigma_{i}(X) \cup \xi^{j-i+1}\, \sigma_{i-1}(X) \cup \cdots \cup \xi^{j}\, \sigma_{0}(X) & \text{si } j > i
  \end{cases}
\]
Ainsi
\[
  \begin{aligned}
    D_{i}(M^{+}(X)) / D_{i-1}(M^{+}(X)) &\simeq \mathbb{Z}^{D_{i}(\Sigma^{+}(X)) - D_{i-1}(\Sigma^{+}(X))} \\
    &\simeq \mathbb{Z}^{\sigma_{i}(X) + \xi \sigma_{i-2}(X) + \xi^{2} \sigma_{i-4}(X) + \cdots}
  \end{aligned}
\]
\struck{$\mathbb{Z}^{\sigma_{i}(X)}$}

\emph{N.B.} \ill{} on voit \ill{} la structure \uncertain{additive} des deux
filtrations, et de l'opération $\xi \cdot$, $M^{+}(X)$ et $M(X)$ sont
complètement déterminés par la connaissance \ill{} la famille
$(\sigma_{i}(X))_{i \geqslant 0}$ des \uncertain{ens.} $\sigma_{i}(X)$.

Pour la structure multiplicative il faut connaître de plus les applications
\[
  \sigma_{i}(X) \times \sigma_{j}(X) \longrightarrow
  \mathbb{Z}^{\left(\coprod_{k \leqslant i+j-2} \sigma_{k}(X)\right)}
\]
\marginal{$x, y \mapsto$ \ill{} définies par \ill{}}
\note{l'indice sous le coproduit se lit mal}

\page{40}
Nous allons expliciter plus bas la détermination de $\sigma^{0}(X)$ et
$\sigma^{1}(X)$.

Il y a lieu également de préciser les opérations $\vee$ et $\lambda^{i}$. Pour
la première, notons que
\[
  \struck{\mathbb{Z}}\mathbb{Q}(-i)\, [\mathcal{M}^{+}_{i}(X)]^{\vee} \subset \mathcal{M}^{+}_{i}(X)
\]
d'où on conclut \ill{} que $x \mapsto \check{x}\, \xi^{i}$ définit \ill{} une
involution. Mais en fait il \uncertain{résulterait} de la théorie \ill{} que
\emph{cette dernière est l'identité}, i.e.\ \emph{l'opération $\vee$ dans
$M_{i}(X)$} \struck{\ill{}} \emph{n'est autre que} $x \mapsto \xi^{-i} x$,
\struck{\ill{}}

Quant aux opérations $\lambda^{i}$, leur détermination devrait se faire en
termes de groupes finis d'opérateurs (tels les $\mathfrak{S}_{p}$ —).
\note{le paragraphe 10) s'arrête ici ; le bas du feuillet est vierge}

\end{document}
